% =============================================================================
% Lecture 17: Responsible Physical AI
% Physical AI: Machine Learning Systems That Sense and Act (Volume IV), Chapter 17
% =============================================================================
% Thesis: authority is fit to grant only when every decision behind an action
% has a named owner, every human role can be performed, and every exposed
% group is counted, in meters and joules where the risk is physical.
% Hook and synthesis: the Cruise collision of 2 October 2023 and the record
% given to regulators afterward (Responsible Physical AI, Sec. 17.1, 17.4.2).
% Running case: the warehouse mobile manipulator's admitted aisle command,
% 1.3 m/s toward a rack end (1.10 m clear distance, 368 kg loaded, 2 m/s^2).
%
% Build: make deck V=4 D=17_responsibility   (figures: images/make_figs.py)
% Numbers: images/calcs.py prints every derived value on these slides.
% Plan:  lecture_map.yml in this folder.
% =============================================================================

\documentclass[aspectratio=169, 12pt]{beamer}
\usepackage{../../assets/beamerthememlsys}
\mlsyssetup{
  volume={Volume IV},
  chapter={Chapter 17},
  chaptertitle={Responsible Physical AI},
  layer=governance,
  logo={../../assets/img/logo-mlsysbook.png},
  instlogo={../../assets/img/logo-harvard.png},
}
\usepackage[T1]{fontenc}
\usepackage[scaled=0.9]{helvet}
\usepackage{courier}
\renewcommand{\familydefault}{\sfdefault}
\usepackage{booktabs,amsmath,array}
\newcolumntype{L}[1]{>{\raggedright\arraybackslash}p{#1}}
\graphicspath{{images/}}
% Volume IV figure roles (volumes/vol4/palette.md), for text beside figures
\definecolor{figperm}{HTML}{047857}
\definecolor{figveto}{HTML}{991B1B}
\definecolor{figseal}{HTML}{92400E}
\setcounter{mlsystotalsections}{5}

\title{Responsible Physical AI}
\author{Vijay Janapa Reddi}
\institute{Harvard University}
\date{}

% Chapter notation, as the book prints it (books/shared/_includes/_macros.tex)
\newcommand{\ArrivalRate}{\lambda_{\text{arr}}}
\newcommand{\IntervalMean}{\bar{T}}
\newcommand{\ServingUtilization}{\rho_{\text{serv}}}
\newcommand{\QueueWaitTime}{W_{q}}
\newcommand{\DetectorMissProbability}{p_{\text{miss}}}
\newcommand{\PerFrameDetectionProbability}{p_{\text{frame}}}
\newcommand{\ReliableDetectionFrameCount}{N_{\text{frames}}}
\newcommand{\DetectionRangeShortfall}{\Delta d_{\text{detect}}}
\newcommand{\CameraFramePeriod}{T_{\text{sample}}}
\newcommand{\SpeedAtDetection}{v}
\newcommand{\CredibleBrakingDeceleration}{a_{\text{brake}}}
\newcommand{\ImpactSpeed}{v_{\text{impact}}}
\newcommand{\KineticEnergy}{E_k}
\newcommand{\Mass}{m}
% A large answer choice for predict and concept-question slides
\newcommand{\choice}[2]{\textbf{\textcolor{mlsysaccent}{#1}}\enspace #2}
% A choice that is highlighted from overlay #1 on (the answer)
\newcommand{\answerat}[3]{\alt<#1->{\colorbox{mlsysaccent!15}{\choice{#2}{#3}}}{\choice{#2}{#3}}}
% A goal or takeaway line with its thumbnail
\newcommand{\thumbline}[2]{%
  \begin{minipage}[c]{2.6cm}\centering\includegraphics[width=2.5cm,height=0.95cm,keepaspectratio]{#1}\end{minipage}%
  \hspace{0.3cm}\begin{minipage}[c]{0.74\textwidth}\small #2\end{minipage}}
% One line of a worked example: shown from overlay #1 on
\newcommand{\wstep}[2]{\uncover<#1->{#2}\\[0.2cm]}
% An exact excerpt of a real document, at a legible size
\newenvironment{docblock}{\par\footnotesize\raggedright\setlength{\parskip}{3pt}}{\par}

\begin{document}

% =============================================================================
% TITLE AND HOOK
% =============================================================================
\mlsysnote{[0.5 min] [CORE] Read the thesis aloud once. Lecture 16 ended with a signed release manifest. Today asks who stands behind the decisions that manifest does not name. The notes total about 71 minutes; the core minutes total about 58, and cutting every [OPTIONAL] frame leaves a complete lecture. Cut first the published-miss-rate conversion, the legible-intent artifact, the penalties chart, and the information-and-authority frame; then the stage-hours, frame-count, and exposure frames.}
\mlsystitle{Responsible Physical AI}{Authority is fit to grant only when each decision behind an action has an owner, each human role can be performed, and each exposed group is counted in meters and joules.}{cover_responsibility_blueprint.png}

\begin{frame}{Who answered for the Cruise collision?}
  \mlsysactivity{Vote}{2 min}
  \begin{columns}[c]
    \begin{column}{0.5\textwidth}
      \centering\safeimg[width=\textwidth,height=3.9cm,keepaspectratio]{cruise-robotaxi.jpg}\\[0.05cm]
      {\footnotesize A Cruise robotaxi in San Francisco, 2022}
    \end{column}
    \begin{column}{0.48\textwidth}
      \small San Francisco, 2 October 2023. A driverless car ran over a pedestrian, then pulled toward the curb with her beneath it.\\[0.25cm]
      {\normalsize Who answered for the pullover?}\\[0.15cm]
      {\small\choice{A}{the classifier engineers}\\
      \choice{B}{remote assistance}\\
      \choice{C}{the executives}\\
      \answerat{2}{D}{no one named for the whole}}
    \end{column}
  \end{columns}
  \mlsyscredit{Photo: Dietmar Rabich, CC BY-SA 4.0, Wikimedia Commons (not the vehicle in the collision). Responsible Physical AI, \S17.1.}
  \note{[3 min] [CORE] NARRATE: a human-driven car struck a pedestrian and threw her into the path of a driverless Cruise vehicle, which braked hard, ran over her, and stopped. With her still beneath it, the vehicle began a pullover and moved about 6.1 m (20 ft) at up to 7.7 mph before it stopped (Quinn Emanuel 2024, as cited in \S17.1). ENGAGE: votes by hand, 30 seconds. Step 2: D. A post-collision classifier placed her at the side of the vehicle, and the fallback it chose moved the car without evidence that the space beneath it was clear; Lecture 12 owns that mechanism. Each group held one piece truthfully: classifier and fallback engineers, remote assistance, the report writers, the executives. No record said who had decided that a driverless car should pull over after a side impact without checking beneath itself. WARN: students who pick A treat a design decision as a coding bug. LINK: today's question is the driving question of the chapter. FLEX: [CORE].}
\end{frame}

\begin{frame}{Regulators penalized Cruise for the record as well as the motion}
  \centering
  \mlsyssteps[width=\textwidth,height=5.0cm,keepaspectratio]{cruise_timeline_step}{2}
  \mlsyscredit{Responsible Physical AI, \S17.1 and \S17.4.2 (CPUC D.24-07-004; NHTSA 2024; U.S. Attorney N.D. Cal.\ 2024; California DMV 2023). CPUC fine derived: 15 days $\times$ \$7,500.}
  \note{[3 min] [CORE] Step 1: the logs and video were good enough for engineers to rebuild the pullover in detail. What failed was the account. On 3 October a Cruise representative told a CPUC analyst by telephone that the car had stopped on impact; the description left out the pullover and the dragging. The CPUC counted 15 days, 3 through 18 October, without a full account, and received the full video on 19 October. On 24 October the DMV suspended the driverless permits. Step 2: the 2024 penalties were for the record. CPUC: 15 days at its \$7,500 daily maximum, \$112,500 (calcs.py). NHTSA: a \$1.5 million consent order for crash reports that omitted the dragging. Justice Department: a \$500,000 criminal fine for a false report. ENGAGE: which of these penalties could a better classifier have avoided? (Expected: none of them; all three concern what was reported.) WARN: the company's own review found no evidence of intent to mislead; the lesson is that disclosure had no owner. FLEX: [CORE].}
\end{frame}

\begin{frame}{Recall: what does a release manifest bind?}
  \mlsysactivity{Retrieval}{1 min}
  \begin{columns}[c]
    \begin{column}{0.56\textwidth}
      {\Large\bfseries Close your notes.}\\[0.2cm]
      {\normalsize Name three things a signed release manifest holds.}\\[0.3cm]
      \uncover<2->{\small\textbf{\textcolor{mlsysaccent}{Verdict:}} which claims the evidence supports\\
      \textbf{\textcolor{mlsysaccent}{Records:}} bound by digest (a hash of each file)\\
      \textbf{\textcolor{mlsysaccent}{Standing conditions:}} what withdraws the grant}
    \end{column}
    \begin{column}{0.4\textwidth}
      \centering\safeimg[width=\textwidth,height=1.7cm,keepaspectratio]{thumb_record_chain.pdf}\\[0.05cm]
      {\footnotesize The record chain: manifest, then today's record, then the register}
    \end{column}
  \end{columns}
  \mlsyscredit{Lecture 16 (release manifest, standing conditions). Responsible Physical AI, \S17.1.}
  \note{[1 min] [CORE] Retrieval, 40 seconds. Step 2: Lecture 16's manifest holds a verdict, the records it binds by digest, and the standing conditions that withdraw it, signed by the lead systems adjudicator. LINK: the manifest names one person. It does not name who decided the aisle speed the warehouse accepts, who decided one remote assistant serves every machine, or who reports an incident, to whom, and how fast. Those are today's decisions. FLEX: [CORE].}
\end{frame}

\begin{frame}{Who answers for a machine that moves on its own?}
  \thumbline{thumb_decision_chain_step3.pdf}{Trace one admitted command back to the six decisions behind it, and find the ones no record holds.}\\[0.16cm]
  \thumbline{thumb_queue_step2.pdf}{Test a human role for time, information, and authority, including the queue one person serves.}\\[0.16cm]
  \thumbline{thumb_ladder_step2.pdf}{Specify a permission stage: what it allows, what it must report, and what withdraws it.}\\[0.16cm]
  \thumbline{thumb_shortfall_step3.pdf}{Convert a per-group detection gap into meters and joules, and size the sample that resolves it.}\\[0.1cm]
  {\footnotesize\textcolor{midgray}{Part IV's law: permission must be earned. Today's record: the approval record, which names an owner for each decision.}}
  \note{[1 min] [CORE] The driving question is the chapter's: who answers for a machine that moves on its own when its motion harms someone? Four goals, one per teaching segment; the fifth segment, records about people and the labor the machine moves, closes on the approval record that collects all of them. LINK: the fourth law, permission must be earned (Lecture 1), lets a machine operate only as far as its evidence supports; today applies it to the decisions behind that evidence. The record chain gains its fifteenth record, the approval record.}
\end{frame}

% =============================================================================
% SECTION: DECISIONS
% =============================================================================
\mlsysnote{[0.5 min] [CORE] First segment: one admitted command, and the people behind it.}
\mlsysdivider{Whose decisions stand behind one admitted command?}{Responsible Physical AI, \S17.1 and \S17.2}
\section{Decisions}

\begin{frame}{The running machine admits 1.3 m/s toward a rack end}
  \begin{columns}[c]
    \begin{column}{0.44\textwidth}
      \centering\safeimg[width=\textwidth,height=4.1cm,keepaspectratio]{mobile-manipulator.jpg}\\[0.05cm]
      {\footnotesize A real mobile manipulator: wheeled base, six-joint arm}
    \end{column}
    \begin{column}{0.54\textwidth}
      \small
      \begin{tabular}{@{}ll@{}}\toprule
        Admitted speed & 1.3 m/s (chosen)\\
        Clear distance & 1.10 m (illustrative)\\
        Loaded mass & 368 kg (derived)\\
        Kinetic energy & 311 J (derived)\\\bottomrule
      \end{tabular}\\[0.3cm]
      The enforcement record shows that this command passed the stopping check at one tick.
    \end{column}
  \end{columns}
  \mlsyscredit{Photo: ER-FLEX (MiR250 base, UR16e arm), Auledas, CC BY 4.0, Wikimedia Commons; a real machine like the book's reference machine, not that machine. Values: MLSysIM registry; energy derived, calcs.py.}
  \note{[1 min] [CORE] LINK: the warehouse mobile manipulator from every lecture of this course. NARRATE: the photo is a real commercial mobile manipulator of the same class. The command we follow is a base velocity of 1.3 m/s toward a rack end where a person may step out, 1.10 m ahead. The machine weighs 368 kg loaded, so it carries one half m v squared, 0.5 times 368 times 1.3 squared, 311 J (calcs.py, derived). ENGAGE: whose decision was the 1.3? (Expected: students say the enforcer's; the enforcer checked it, but the lead systems adjudicator signed it into the release envelope, Lecture 16.) FLEX: [CORE].}
\end{frame}

\begin{frame}{Six decisions by six parties stand behind that one command}
  \centering
  \mlsyssteps[width=\textwidth,height=5.2cm,keepaspectratio]{decision_chain_step}{3}
  \mlsyscredit{Responsible Physical AI, \S17.2, figure 1 (redrawn).}
  \note{[3.5 min] [CORE] Step 1: four decisions leave a trace in records built earlier in the course: the policy weights in the policy manifest (Lecture 6), the enforcer and its thresholds in the enforcement record (Lecture 12), the authority state at that tick in the authority record (Lecture 14), and the operating envelope, aisle speed and floor-friction premise included, in the release manifest (Lecture 16). Step 2: two govern the site rather than the machine: facility management's crossing rule, which lets people enter the aisle only at marked points, and operations' choice of how many machines one assistant serves. No record holds either. Step 3: the approval record collects all six, each with an owner. ENGAGE: which party can truthfully say the harm came from someone else's part? (Expected: every one of them; Thompson 1980 calls this the problem of many hands.) WARN: a learned proposer tempts a seventh excuse, that the model did it. The model proposes and never admits (Lecture 1, the third law), so every admitted action passed a check someone designed under an envelope someone signed. FLEX: [CORE].}
\end{frame}

\begin{frame}{Estimate: how often does a six-link chain resolve?}
  \mlsysactivity{Estimate}{2 min}
  \begin{columns}[c]
    \begin{column}{0.42\textwidth}
      \begin{mlsyscard}{mlsysaccent}\footnotesize
        Each of the 6 decisions can be retrieved, when an investigator asks, with probability $p$, independently.
      \end{mlsyscard}
      \vspace{0.15cm}
      {\small At $p = 0.90$, what share of incidents yields the whole chain?}\\[0.2cm]
      \uncover<2->{\small $0.90^6 = 0.53$, so about half of incidents yield the whole chain.}
    \end{column}
    \begin{column}{0.56\textwidth}
      \centering\mlsyssteps[width=\textwidth,height=5.0cm,keepaspectratio]{chain_prob_step}{2}
    \end{column}
  \end{columns}
  \mlsyscredit{Responsible Physical AI, \S17.2.1 (retrieval probabilities illustrative). Chain values derived, calcs.py.}
  \note{[2.5 min] [CORE] Give the class one minute. Step 2: the whole chain resolves with probability p to the sixth: 0.94 at 0.99 per link, 0.74 at 0.95, and 0.53 at 0.90 (calcs.py). The chapter calls 0.90 roughly what a decision achieves when it is held only in memory and email after its author changes jobs. ENGAGE: why does the curve fall faster than the one-link line? (Expected: a conjunction is weaker than any one of its premises; the same arithmetic made a release argument weaker than its weakest leg in Lecture 16.) WARN: most of these decisions never reach a runtime log. They happen in meetings weeks before a command crosses the causal boundary (the envelope sign-off, the staffing ratio, the permit application) or days after it (what to tell regulators). In the Cruise case, the pullover rule was a design decision made long before 2 October. The remedy is to write each link down as part of the design. FLEX: [CORE].}
\end{frame}

% =============================================================================
% SECTION: OVERSIGHT
% =============================================================================
\mlsysnote{[0.5 min] [CORE] Second segment: the people the safety argument relies on, and whether they can do what it asks.}
\mlsysdivider{When does a person in the loop add protection?}{Responsible Physical AI, \S17.3}
\section{Oversight}

\begin{frame}{Oversight needs time, information, and authority}
  \begin{columns}[c]
    \begin{column}{0.46\textwidth}
      \centering\safeimg[width=\textwidth,height=3.6cm,keepaspectratio]{uber-test-vehicle.jpg}\\[0.05cm]
      {\footnotesize An Uber self-driving test vehicle. In Tempe in 2018, its backup operator answered in court.}
    \end{column}
    \begin{column}{0.52\textwidth}
      \small A role is oversight a person can perform only if its holder has\\[0.12cm]
      \textbf{\textcolor{mlsysaccent}{time}} to act before the hazard arrives,\\
      \textbf{\textcolor{mlsysaccent}{information}} about what the machine is doing, and\\
      \textbf{\textcolor{mlsysaccent}{authority}} to change it.\\[0.2cm]
      A role that fails any of the three moves blame onto its holder.
    \end{column}
  \end{columns}
  \mlsyscredit{Photo: foam from Elgaland Vargaland, CC BY-SA 2.0, Wikimedia Commons (not the Tempe vehicle). Maricopa County Attorney 2023; Elish 2019; Responsible Physical AI, \S17.3, Definition 17.1.}
  \note{[2.5 min] [CORE] NARRATE: a human oversight role is a duty the design assigns to a person, such as stopping the machine, taking over a channel, or approving a motion, that the safety argument relies on. Read the three conditions. Santoni de Sio and van den Hoven (2018) state the same requirement for meaningful human control. LINK: after the Uber test vehicle killed a pedestrian in Tempe in 2018 (Lecture 1), the person who answered in court was the backup operator, who pleaded guilty to endangerment in 2023. Elish (2019) calls a person in that position a moral crumple zone: like a car's crumple zone, the person absorbs the impact, except the impact is blame, and what it protects is the system's reputation. ENGAGE: name a role in the running machine's design that fails one of the three. (Expected: the remote assistant for a person stepping out; next slide.) WARN: adding a person to the loop does not make a design responsible. FLEX: [CORE].}
\end{frame}

\begin{frame}{Can the assistant stop the rack-end hazard?}
  \mlsysactivity{Diagnose}{2 min}
  \centering
  \safeimg[width=\textwidth,height=3.3cm,keepaspectratio]{time_test.pdf}\\[0.15cm]
  \uncover<2->{\small The takeover needs about 5$\times$ the time the hazard leaves. The role is \textbf{\textcolor{figveto}{INFEASIBLE}} as a countermeasure.}
  \mlsyscredit{Responsible Physical AI, \S17.3.1. Out-of-loop takeover: MLSysIM registry, from Lecture 14. Window, travel, and ratio derived, calcs.py.}
  \note{[2.5 min] [CORE] Setup: a person steps out 1.10 m ahead of the machine at 1.3 m/s; the rack end arrives in 1.10 over 1.3, 0.85 s. A remote assistant who has been monitoring rather than driving needs the design out-of-loop takeover, 4.47 s (Lecture 14 derives it from measured take-over times plus video and handshake delays), during which the machine travels 1.3 times 4.47, 5.8 m. Step 2: 4.47 over 0.85 is 5.3, about five times (calcs.py). ENGAGE: can training or diligence close the gap? (Expected: no; even the in-loop takeover of 0.605 s leaves no time to brake, Lecture 14.) WARN: a design that names the assistant as this countermeasure has recorded in advance who will be blamed. The book's design stops the machine first and asks for help from standstill, where waiting costs time and no clearance. The approval record lists, per role, the response time required, the time the hazard leaves, and a verdict: FEASIBLE, INFEASIBLE, or UNKNOWN. FLEX: [CORE].}
\end{frame}

\begin{frame}{One assistant for twenty machines is busy half the time}
  \begin{columns}[c]
    \begin{column}{0.4\textwidth}
      \small Utilization:\\[0.05cm]
      \centering$\ServingUtilization = \ArrivalRate\,\IntervalMean$\\[0.15cm]
      \raggedright Mean wait (random arrivals and service):\\[0.05cm]
      \centering$\QueueWaitTime = \dfrac{\ServingUtilization\,\IntervalMean}{1-\ServingUtilization}$\\[0.2cm]
      \raggedright\footnotesize 0.5 requests per machine-hour, 3 min each (illustrative). At 20 machines, 1 request in 100 waits over 23 min.
    \end{column}
    \begin{column}{0.58\textwidth}
      \centering\mlsyssteps[width=\textwidth,height=5.2cm,keepaspectratio]{queue_step}{2}
    \end{column}
  \end{columns}
  \mlsyscredit{Responsible Physical AI, Napkin Math 17.1 (inputs illustrative); M/M/1 queue. Utilization, waits, and the 99th percentile derived, calcs.py.}
  \note{[3.5 min] [CORE] NARRATE: an assistant who serves many machines is a server in a queue. $\ArrivalRate$ is the fleet's request rate and $\IntervalMean$ the mean time a request holds the assistant. Step 1: 10 machines send 5 requests an hour at 180 s each, utilization 0.25, mean wait 60 s; 20 machines give 0.50 and 180 s. Step 2: 32 machines give 0.80 and 720 s, 12 minutes; at 20 machines the 99th-percentile wait is 1,408 s, 23.5 min (calcs.py). ENGAGE: Checkpoint 17.1 has 24 machines at 0.8 requests per machine-hour and 75 s per request. Utilization? (Expected: 0.40, mean wait 50 s; still far longer than 0.85 s.) WARN: every one of these waits exceeds the rack-end window, so the machines-per-assistant ratio is a productivity budget and never a safety layer. A cap at utilization 0.5 limits this assistant to 20 machines, and the approval record lists that cap as the decision of whoever set it. LINK: Lecture 14 admits a remote-assistance session only in a state where waiting costs no clearance. FLEX: [CORE].}
\end{frame}

\begin{frame}{Authority means a stop channel the permission path honors}
  \begin{columns}[c]
    \begin{column}{0.36\textwidth}
      \centering\safeimg[width=\textwidth,height=4.6cm,keepaspectratio]{agv-stop-button.jpg}\\[0.05cm]
      {\footnotesize Red stop button on the mast; safety scanner at the floor}
    \end{column}
    \begin{column}{0.62\textwidth}
      \footnotesize
      \setlength{\tabcolsep}{4pt}\begin{tabular}{@{}L{3.0cm}L{2.0cm}L{1.5cm}L{1.1cm}@{}}\toprule
        \textbf{Role and hazard} & \textbf{Time} & \textbf{Sees} & \textbf{Acts}\\\midrule
        Assistant, stalled machine & FEASIBLE & console & stop\\
        Assistant, rack end & \textcolor{figveto}{INFEASIBLE} & console & stop\\
        Fleet supervisor & UNKNOWN & dashboard & \textcolor{figveto}{report}\\\bottomrule
      \end{tabular}\\[0.25cm]
      \small One person can hold two duties with two verdicts.
    \end{column}
  \end{columns}
  \mlsyscredit{Photo: Marc Auledas Noguera, CC BY 4.0, Wikimedia Commons (cropped). Responsible Physical AI, \S17.3.3, Checkpoint 17.1.}
  \note{[1.5 min] [OPTIONAL] NARRATE: the photo shows a real authority channel: a stop button a nearby person can press, wired to stop the truck without asking anyone. The test for authority is mechanical: the person must be able to stop the machine through a channel the permission path honors within the time the hazard leaves. A role whose only output is a report, a flag raised to a supervisor or a note in a log, has no authority over the motion. Information matters too: a console that shows video but not which channel the machine still owns leaves the operator guessing (mode confusion, Lecture 14). The table fills the human-roles block for three people from Checkpoint 17.1. ENGAGE: why does the remote assistant appear twice? (Expected: the same person passes all three tests for a stranded machine at standstill and fails the time test for a person stepping out; both facts belong in the record.) FLEX: [OPTIONAL].}
\end{frame}

% =============================================================================
% SECTION: PERMISSION
% =============================================================================
\mlsysnote{[0.5 min] [CORE] Third segment: who may be exposed to the machine, and when.}
\mlsysdivider{How is permission to operate granted, and how is it taken back?}{Responsible Physical AI, \S17.4}
\section{Permission}

\begin{frame}{Permission to operate is granted one stage at a time}
  \begin{columns}[c]
    \begin{column}{0.7\textwidth}
      \centering\mlsyssteps[width=\textwidth,height=4.9cm,keepaspectratio]{ladder_step}{2}
    \end{column}
    \begin{column}{0.28\textwidth}
      \centering\safeimg[width=\textwidth,height=2.6cm,keepaspectratio]{waymo-driverless.jpg}\\[0.05cm]
      {\footnotesize A Waymo car with no one in the driver's seat, San Francisco, 2023}
    \end{column}
  \end{columns}
  \mlsyscredit{Photo: Dllu, CC BY 4.0, Wikimedia Commons. California Code of Regulations, Title 13, Articles 3.7 and 3.8; CPUC 2023; Responsible Physical AI, \S17.4, Definition 17.2.}
  \note{[3 min] [CORE] NARRATE: a permission stage is a grant, made by the party entitled to grant it, that lets a named configuration operate in a named domain, among named groups of people, on named evidence, under named reporting duties. Step 1: California separates testing with a safety driver, testing without one, and deployment, each its own permit. The CPUC adds rungs for passengers: before August 2023 Cruise could carry paying passengers without a safety driver only in parts of San Francisco between 10 p.m. and 6 a.m.; on 10 August 2023 the commission removed the time limit. Step 2: a warehouse has the same ladder: commissioning with a spotter, unattended running in a zone people cannot enter, then the open floor. The facility owner, the integrator, and workers' representatives grant these. ENGAGE: how does a stage differ from a release verdict? (Expected: the verdict says which claims the evidence supports, Lecture 16; the stage says who may be exposed and what must be reported.) FLEX: [CORE].}
\end{frame}

\begin{frame}{Ten machines need 19 days to leave the supervised stage}
  \begin{columns}[c]
    \begin{column}{0.4\textwidth}
      \small Claim: fewer than 1 safety-relevant intervention per 1,000 h, at 95\% confidence, after observing none.\\[0.2cm]
      Hours needed $= \ln(1/0.05) \times 1{,}000$ h.\\[0.2cm]
      \footnotesize A spotter masks failures, so each intervention is reconstructed.
    \end{column}
    \begin{column}{0.58\textwidth}
      \centering\safeimg[width=\textwidth,height=5.2cm,keepaspectratio]{stage_hours.pdf}
    \end{column}
  \end{columns}
  \mlsyscredit{Responsible Physical AI, \S17.4.1 (fleet size and hours illustrative). Hours and days derived, calcs.py.}
  \note{[2 min] [OPTIONAL] LINK: this is Lecture 7's exposure wall applied to a stage exit. NARRATE: to claim fewer than one safety-relevant intervention per 1,000 operating hours at 95 percent confidence after observing none, the fleet needs ln of 20 over 10 to the minus 3 per hour, 2,996 failure-free hours. Ten machines running 16 hours a day accumulate 160 hours a day and reach it on day 18.7, so 19 days (calcs.py). WARN: the protection the stage provides hides its own evidence. When a spotter or safety driver intervenes, the machine's error never reaches its consequence, and the stage records an intervention. Counting failures means reconstructing, for each intervention, whether the machine would have failed without it, which is why disengagement reports must say why control was transferred. ENGAGE: what resets the count? (Next slide.) FLEX: [OPTIONAL].}
\end{frame}

\begin{frame}{True or false: an update keeps the stage it ran under?}
  \mlsysactivity{Vote}{1.5 min}
  \centering
  {\normalsize ``Our fleet earned unattended operation last month. This week's policy update runs under the same stage.''}\\[0.4cm]
  {\Large\choice{A}{True}\qquad\answerat{2}{B}{False}}\\[0.4cm]
  \uncover<2->{\small The 2,996 hours belong to the configuration that ran them. A new policy, enforcer build, or site layout starts its own count.}
  \mlsyscredit{Responsible Physical AI, \S17.4.1 and \S17.10 (pitfall: carrying a permission stage across a configuration change).}
  \note{[1.5 min] [CORE] Vote, then step 2: false. A stage authorizes a configuration in a domain under reporting duties. A new policy, a new enforcer build, or an expanded operating domain is a different configuration, and the failure-free hours behind the old stage do not transfer. Lecture 16 made the same point for release claims: certification authorizes a configuration. The countermeasure is mechanical: bind each stage to a configuration digest in the permit block of the approval record, so an update reopens the stage along with the release claims it touches. WARN: students who answer true think of a permit as a company license. It belongs to a configuration and carries duties that, if broken, withdraw it. FLEX: [CORE].}
\end{frame}

\begin{frame}[t]{The DMV cited unsafe motion and misrepresented information}
  \begin{columns}[T]
    \begin{column}{0.62\textwidth}
      \begin{docblock}
        \textcolor{midgray}{DMV Statement on Cruise LLC Suspension, 24 October 2023}\\[0.08cm]
        ``13 CCR \S228.20 (b) (6) -- Based upon the performance of the vehicles, the Department determines the manufacturer's vehicles are not safe for the public's operation.''\\[0.06cm]
        ``13 CCR \S228.20 (b) (3) -- The manufacturer has misrepresented any information related to safety of the autonomous technology of its vehicles.''
      \end{docblock}
    \end{column}
    \begin{column}{0.36\textwidth}
      \small Withdrawal steps the company down the ladder.\\[0.3cm]
      One ground concerns the \textbf{\textcolor{figveto}{motion}}; the other concerns the \textbf{\textcolor{figseal}{record}}.
    \end{column}
  \end{columns}
  \mlsyscredit{California DMV, dmv.ca.gov, ``DMV Statement on Cruise LLC Suspension,'' 24 Oct.\ 2023 (two of four grounds quoted). Responsible Physical AI, \S17.4.2.}
  \note{[2 min] [CORE] NARRATE: an exact excerpt of the public statement; it cites four grounds, and two are quoted. The suspension removed driverless testing and deployment and left the safety-driver permit in place, so permission was withdrawn the way it was granted, one stage at a time. The first ground is a finding about motion; the second is a finding about the record. LINK: the incident disclosure duty is the obligation, attached to a permit, to report an incident completely and on time to the parties the permit names. California requires a collision report within ten days and an annual disengagement report; NHTSA's order, in its 2023 version, required a crash involving a death, a hospital-treated injury, or a vulnerable road user to be reported within one day, with an update at ten. ENGAGE: who should own the first report after a collision? (Expected: a named role in the permit block, with author, recipient, deadline, and evidence attached.) FLEX: [CORE].}
\end{frame}

\begin{frame}{Three penalties for the record totaled \$2.1 million}
  \begin{columns}[c]
    \begin{column}{0.58\textwidth}
      \centering\safeimg[width=\textwidth,height=4.6cm,keepaspectratio]{fines.pdf}
    \end{column}
    \begin{column}{0.4\textwidth}
      \small The permit block of the approval record keeps a disclosure log: each report, its author, its recipient, its deadline, and the evidence attached.
    \end{column}
  \end{columns}
  \mlsyscredit{Responsible Physical AI, \S17.4.2 (CPUC D.24-07-004; NHTSA consent order, Sept.\ 2024; U.S. Attorney N.D. Cal., Nov.\ 2024). Total derived, calcs.py.}
  \note{[1.5 min] [OPTIONAL] NARRATE: \$112,500 from the CPUC for 15 days at its daily maximum; \$1.5 million from NHTSA for crash reports that omitted the dragging; \$500,000 as a criminal fine under a deferred prosecution agreement for a false report meant to influence the federal investigation. Total \$2,112,500 (calcs.py). ENGAGE: what is the engineering half of the disclosure duty? (Expected: a release manifest that binds raw traces to the configuration that produced them, Lecture 16, so an assessor can open them; the disclosure log is the other half, showing what was actually sent, to whom, and when.) WARN: the dollar amounts are small next to the company's losses; the lesson is that the permit holder's record is itself a premise of its authority. FLEX: [OPTIONAL].}
\end{frame}

% =============================================================================
% SECTION: DETECTION
% =============================================================================
\mlsysnote{[0.5 min] [CORE] Fourth segment: who the machine sees well, and what a gap costs in meters and joules.}
\mlsysdivider{What does a detection gap between groups of people cost?}{Responsible Physical AI, \S17.5 and \S17.6}
\section{Detection}

\begin{frame}{Published pedestrian detectors miss some groups far more often}
  \centering
  \safeimg[width=\textwidth,height=5.2cm,keepaspectratio]{detect_data.pdf}\\[0.05cm]
  \small Children: 43.89\% missed, adults 23.75\% (averages). Darker skin: AP75 55.9\%, lighter 67.1\%.
  \mlsyscredit{Li et al., ``Bias Behind the Wheel,'' ACM TOSEM 34(3), 2025, Table 7 (eight detectors). Wilson, Hoffman, Morgenstern, arXiv 1902.11097, 2019, Table 4 (COCO-trained models on BDD100K).}
  \note{[3 min] [CORE] NARRATE left: Li and colleagues labeled pedestrians in four driving datasets and tested eight widely used detectors. Every detector missed children more often; the average gap is 20.14 percentage points (calcs.py recomputes it from Table 7). Right: Wilson, Hoffman, and Morgenstern labeled pedestrians in BDD100K by Fitzpatrick skin type. AP75, average precision under a strict box-overlap criterion, is lower for darker skin in all eight models, 55.9 against 67.1 percent on average; time of day and occlusion did not explain it. ENGAGE: Li et al. found almost no skin-tone gap in miss rate on their data. Does that contradict Wilson? (Expected: no; different datasets, a different metric, and box size all differ. A per-group rate must be measured for the machine's own detector and site.) WARN: these are road datasets; a warehouse has its own groups: kneeling workers, people carrying cartons, wheelchair users, workers without a high-visibility vest. FLEX: [CORE].}
\end{frame}

\begin{frame}{A group seen less often needs more frames before detection counts}
  \begin{columns}[c]
    \begin{column}{0.44\textwidth}
      \small Frames in view before a person counts as detected:\\[0.1cm]
      \centering$\ReliableDetectionFrameCount = \left\lceil \dfrac{\ln \DetectorMissProbability}{\ln(1-\PerFrameDetectionProbability)} \right\rceil$\\[0.2cm]
      \raggedright\footnotesize $\PerFrameDetectionProbability$: chance of a detection in one frame. $\DetectorMissProbability$: allowed chance of missing every frame, $10^{-3}$. Frames assumed independent.
    \end{column}
    \begin{column}{0.54\textwidth}
      \centering\safeimg[width=\textwidth,height=5.1cm,keepaspectratio]{miss_frames.pdf}
    \end{column}
  \end{columns}
  \mlsyscredit{Responsible Physical AI, \S17.5.1, equation 3 (per-frame rates illustrative; $10^{-3}$ chosen). Values derived, calcs.py.}
  \note{[2 min] [OPTIONAL] NARRATE: the plot shows the chance that a person in view is missed in every one of n consecutive frames, (1 minus p) to the n. A group seen in 95 percent of frames falls below the allowed 1 in 1,000 after 3 frames; a group seen in 75 percent needs 5 (calcs.py). ENGAGE: why does a 20-point gap in per-frame rate become only 2 frames? (Expected: the miss probability falls geometrically, so the frame count grows with the logarithm.) WARN: independence is the best case. A detector that systematically fails on a pose, a kneeling worker for example, misses many frames in a row, and the count grows well beyond this. LINK: this matters only where the design relies on a learned detector to recognize a person (to slow down near people, keep a wider separation, or predict a path). A certified safety scanner that stops for any object above its resolution does not care who the object is. FLEX: [OPTIONAL].}
\end{frame}

\begin{frame}{What reaches a kneeling worker two frames late?}
  \mlsysactivity{Predict}{2 min}
  \begin{columns}[c]
    \begin{column}{0.48\textwidth}
      \begin{mlsyscard}{mlsysaccent}\footnotesize
        Stop sized for standing workers, seen in 95\% of frames (3 frames). Kneeling workers: 75\% (5 frames).\\
        60 Hz camera, 1.3 m/s, 2 m/s$^2$, 368 kg.
      \end{mlsyscard}
      \vspace{0.1cm}
      {\small Contact speed and energy?}\\[0.15cm]
      \uncover<2->{\small
      \begin{tabular}{@{}ll@{}}
        $\DetectionRangeShortfall = \SpeedAtDetection\,\CameraFramePeriod\,\Delta\ReliableDetectionFrameCount$ & 43 mm\\[0.05cm]
        $\ImpactSpeed = \sqrt{2\,\CredibleBrakingDeceleration\,\DetectionRangeShortfall}$ & 0.42 m/s\\[0.05cm]
        $\KineticEnergy = \Mass\,\CredibleBrakingDeceleration\,\DetectionRangeShortfall$ & 32 J
      \end{tabular}}
    \end{column}
    \begin{column}{0.5\textwidth}
      \centering\safeimg[width=\textwidth,height=4.6cm,keepaspectratio]{shortfall_step1.pdf}
    \end{column}
  \end{columns}
  \mlsyscredit{Responsible Physical AI, Napkin Math 17.2 (rates and deceleration illustrative; mass derived; camera datasheet). Values derived, calcs.py.}
  \note{[2.5 min] [CORE] The plot is the stop sized for standing workers: speed against distance from first view, ending at rest just short of the person line. Give the class a minute to guess: zero, a nudge, or a strike? Step 2: two extra frames at 16.7 ms are 33.3 ms; at 1.3 m/s that is 43 mm of travel before braking can begin, well inside the 0.42 m braking distance. Braking at 2 m/s squared over a stop that was sized to end exactly at the standing worker leaves the machine moving when it reaches the kneeling one: square root of 2 times 2 times 0.043, 0.42 m/s, carrying 368 times 2 times 0.043, 32 J (calcs.py). ENGAGE: the energy grows in proportion to which quantity? (Expected: the distance lost; the last form of equation 5.) WARN: students expect a 20-point recall gap to cost a little recall. It costs contact, because the stop had no margin for the slower group. The thumbnail is the next slide's speed-distance plot. FLEX: [CORE].}
\end{frame}

\begin{frame}{One stop ends short of one group and strikes another}
  \begin{columns}[c]
    \begin{column}{0.54\textwidth}
      \centering\mlsyssteps[width=\textwidth,height=5.2cm,keepaspectratio]{shortfall_step}{3}
    \end{column}
    \begin{column}{0.44\textwidth}
      \footnotesize
      \begin{tabular}{@{}rrrr@{}}\toprule
        \textbf{Per frame} & \textbf{Lost} & \textbf{Impact} & \textbf{Energy}\\\midrule
        0.95 & 0 mm & 0 m/s & 0 J\\
        0.85 & 22 mm & 0.29 m/s & 16 J\\
        0.75 & 43 mm & 0.42 m/s & 32 J\\
        0.60 & 108 mm & 0.66 m/s & 80 J\\\bottomrule
      \end{tabular}\\[0.3cm]
      \small Fairness for a moving machine is measured in meters and joules.
    \end{column}
  \end{columns}
  \mlsyscredit{Responsible Physical AI, \S17.5.1, Table 1 and figure 2 (redrawn; illustrative rates). Values derived, calcs.py.}
  \note{[2 min] [CORE] Step 1: speed against distance from the moment a person enters the detector's view. The reference group is detected after 3 frames; after the 133.6 ms pre-brake delay the machine brakes and stops exactly at the person. Step 2: the group seen in 75 percent of frames brakes 43 mm later and arrives at 0.42 m/s. Step 3: a group seen in 60 percent needs 8 frames, 5 extra, loses 108 mm, and is struck at 0.66 m/s with 80 J. The table is the chapter's Table 1 (calcs.py). ENGAGE: compare remedies for the kneeling-worker gap: a lower speed near racks, a sensor that does not depend on pose, or a zone rule that keeps the machine out while people kneel there. What evidence does each need? (Expected: the speed cut needs only the per-group rate; the pose-independent sensor needs its own qualification, Lecture 8's observation contract; the zone rule needs compliance, a premise no sensor checks.) FLEX: [CORE].}
\end{frame}

\begin{frame}{Published miss rates for children cost 65 to 130 mm at aisle speed}
  \centering
  \safeimg[width=\textwidth,height=4.9cm,keepaspectratio]{li_conv.pdf}\\[0.05cm]
  \small Derived: each detector's per-image miss rate treated as an independent per-frame rate.
  \mlsyscredit{Miss rates: Li et al., ACM TOSEM 2025, Table 7. Conversion at the running machine's speed, camera, mass, and deceleration (Responsible Physical AI, equations 3 to 5); derived, calcs.py.}
  \note{[1 min] [OPTIONAL] NARRATE: this frame joins published data to the chapter's equations. For each of Li's eight detectors, take the adult and child miss rates as per-frame miss probabilities, compute the frames each group needs at 1 in 1,000, and convert the difference at 1.3 m/s. Children lose 65 to 130 mm and would be struck at 0.51 to 0.72 m/s with 48 to 96 J; with the averages, 4 extra frames, 87 mm, 0.59 m/s, 64 J (calcs.py). WARN: this is a derived illustration. Li measured misses on single road images, not frame sequences in a warehouse, and real misses are correlated across frames, which makes the true shortfall larger. ENGAGE: which detector would you choose from this plot alone? (Expected: none from this alone; the absolute miss rate matters too, and Faster RCNN has the lowest of both.) FLEX: [OPTIONAL].}
\end{frame}

\begin{frame}{A 200-person group cannot tell 0.95 detection from 0.92}
  \begin{columns}[c]
    \begin{column}{0.56\textwidth}
      \centering\safeimg[width=\textwidth,height=4.6cm,keepaspectratio]{intervals.pdf}
    \end{column}
    \begin{column}{0.42\textwidth}
      \small A 10,000-person test set with a 2\% group, both at 0.95.\\[0.2cm]
      Resolving 0.95 from 0.92 at 80\% power takes \textbf{1,059} people in \textbf{each} group.\\[0.2cm]
      \footnotesize Wilson et al.\ had 387 and 100 large pedestrians in their holdout.
    \end{column}
  \end{columns}
  \mlsyscredit{Responsible Physical AI, \S17.5.2 (set illustrative); exact (Clopper-Pearson) intervals and two-proportion sample size, calcs.py. Wilson et al., arXiv 1902.11097, \S4.2.}
  \note{[2 min] [CORE] NARRATE: same detection rate, 0.95, in both groups. The exact 95 percent interval for 9,800 people runs 0.945 to 0.954; for 200 people, 0.910 to 0.976, about 8 times wider (calcs.py). It cannot tell 0.95 from 0.92. Telling them apart with 80 percent power takes 1,059 people per group, about 11 percent of the whole set for every group the evaluation must resolve. Wilson and colleagues hit this wall: their holdout had 387 and 100 large pedestrians, and they estimated 12,000 and 4,000 to verify a 0.05 gap. ENGAGE: will a group that is 2 percent of the workforce reach 1,059 by chance? (Expected: no; the evaluation must recruit it, so the per-group sample size goes into the plan before collection.) LINK: the observation contract (Lecture 8) states detection per group, the fault list (Lecture 15) seeds kneeling and carrying people, and the evaluation record (Lecture 7) reports each group with its interval. A group with no entry is an unknown premise. FLEX: [CORE].}
\end{frame}

\begin{frame}{People who never chose the machine can carry most of its risk}
  \begin{columns}[c]
    \begin{column}{0.42\textwidth}
      \centering\safeimg[width=\textwidth,height=4.4cm,keepaspectratio]{starship-crosswalk.jpg}\\[0.05cm]
      {\footnotesize A sidewalk delivery robot on a crosswalk, Finland}
    \end{column}
    \begin{column}{0.56\textwidth}
      \centering\safeimg[width=\textwidth,height=4.3cm,keepaspectratio]{exposure.pdf}\\[0.05cm]
      \raggedright\footnotesize 80\% of encounters at three crossings used by 10\% of workers (illustrative).
    \end{column}
  \end{columns}
  \mlsyscredit{Photo: JIP, CC BY-SA 4.0, Wikimedia Commons. Responsible Physical AI, \S17.6 (shares illustrative); ratios derived, calcs.py.}
  \note{[2 min] [OPTIONAL] NARRATE: the robot's operator chose it, and so did the customer waiting for food. The pedestrians on that crosswalk did not. Risk carried by people who did not choose it is unconsented risk, and no aggregate figure reveals it. On the warehouse floor the crossing rule buys the machine its speed by asking pedestrians to wait and comply. If 80 percent of encounters happen at three crossings used by 10 percent of workers, those workers carry 0.80 over 0.10, 8 times the average exposure per person, and everyone else carries 0.20 over 0.90, 22 percent of it (calcs.py). ENGAGE: an operator reports that injuries per machine-hour fell this year. What question do you ask? (Expected: per group and per location, and whether each group chose its exposure.) WARN: routing, operating-domain boundaries, and operating hours are decisions about who bears the risk. LINK: Lecture 16's fleet evidence, reported one outcome at a time, still needs a per-group split. FLEX: [OPTIONAL].}
\end{frame}

\begin{frame}[t]{A wheelchair user's report paused a delivery fleet within two hours}
  \begin{columns}[T]
    \begin{column}{0.58\textwidth}
      \begin{docblock}
        \textcolor{midgray}{The Pitt News, 21 October 2019}\\[0.08cm]
        Ackerman, on Twitter: ``i (in a wheelchair) was just trapped on forbes ave by one of these robots, only days after their independent roll out.''\\[0.06cm]
        Pitt: ``We have paused testing to review the technology further and ensure everyone's safety.''
      \end{docblock}
    \end{column}
    \begin{column}{0.4\textwidth}
      \small People protect themselves by predicting where a machine will go.\\[0.2cm]
      \textbf{Legible} motion reveals the goal early. Signals are shown only where measured detection supports them.
    \end{column}
  \end{columns}
  \mlsyscredit{E.~Wolfe, ``Pitt pauses testing of Starship robots due to safety concerns,'' The Pitt News, 21 Oct.\ 2019 (quotes verbatim). Dragan et al., HRI 2013; Lee and See 2004; Responsible Physical AI, \S17.6.1.}
  \note{[1.5 min] [OPTIONAL] NARRATE: a real accessibility complaint. A doctoral student who uses a wheelchair reported that a delivery robot idling at a curb cut on Forbes Avenue blocked her way onto the sidewalk; Pitt paused testing within hours, and Starship updated its map for that intersection. Starship later said its video showed her passing the robot; the account is disputed, and the fleet's response is not. LINK: Dragan, Lee, and Srinivasa distinguish predictable motion, which matches what an observer who knows the goal expects, from legible motion, which lets an observer who does not know the goal infer it early. A base that drifts toward the side it will turn into is legible. ENGAGE: why show a signal only where detection is measured? (Expected: Lee and See showed reliance follows trust; a signal promising that the machine slows for people invites a kneeling worker to rely on a detector that sees kneeling workers late.) FLEX: [OPTIONAL].}
\end{frame}

% =============================================================================
% SECTION: RECORDS
% =============================================================================
\mlsysnote{[0.5 min] [CORE] Last segment: the records about people, the labor the machine moves, and the record that collects every decision.}
\mlsysdivider{What does the machine record about people, and what work does it create?}{Responsible Physical AI, \S17.7 to \S17.9}
\section{Records}

\begin{frame}[t]{Event-triggered retention keeps 1\% of the footage}
  \begin{columns}[T]
    \begin{column}{0.44\textwidth}
      \begin{docblock}
        \textcolor{midgray}{GDPR, Article 5(1)}\\[0.06cm]
        (c) ``limited to what is necessary''\\[0.04cm]
        (e) ``kept in a form which permits identification of data subjects for no longer than is necessary''
      \end{docblock}
    \end{column}
    \begin{column}{0.54\textwidth}
      \vspace{0.1cm}
      \centering\safeimg[width=\textwidth,height=3.4cm,keepaspectratio]{retention.pdf}\\[0.1cm]
      \raggedright\small 20 events a day, each 5 s before and 25 s after: 30 s kept per event, a 96$\times$ reduction.
    \end{column}
  \end{columns}
  \mlsyscredit{Regulation (EU) 2016/679, Art.\ 5(1)(c) and (e), excerpts. Responsible Physical AI, \S17.7 (event count and day illustrative); share and reduction derived, calcs.py.}
  \note{[3 min] [CORE] NARRATE: every record a release relies on also records people: the authority record names the operator behind each takeover (Lecture 14), the black box holds frames of whoever stood in the aisle, and the demonstration archive binds each episode to its demonstrator (Lecture 5's provenance record). Kept forever and open to anyone, the safety store becomes workplace surveillance. The excerpt is GDPR's minimization and storage rules. The engineering answer: keep the 5 s before a trigger that reconstructing a takeover needs and 25 s after, 30 s per event; at 20 events in a 16-hour day the machine keeps 10 minutes of 960, about 1 percent, a 96-fold reduction (calcs.py). ENGAGE: what turns a declaration into a premise the gate can check? (Expected: a standing condition, Lecture 16, with an access and retention audit as its detector and suspension of the use as its response.) WARN: keeping everything feels safe; records without a declared purpose drift into uses no one approved. FLEX: [CORE].}
\end{frame}

\begin{frame}[t]{Discuss: may a site rank operators by their takeovers?}
  \mlsysactivity{Discuss}{2 min}
  \begin{columns}[T]
    \begin{column}{0.56\textwidth}
      \begin{docblock}
        \textcolor{midgray}{EU AI Act, Annex III, point 4(b): high-risk use}\\[0.04cm]
        ``AI systems intended to be used \ldots{} to monitor and evaluate the performance \ldots{} of persons in such relationships.''\\[0.08cm]
        \textcolor{midgray}{Article 26(7)}\\[0.04cm]
        ``\ldots{} deployers who are employers shall inform workers' representatives and the affected workers that they will be subject to the use of the high-risk AI system.''
      \end{docblock}
    \end{column}
    \begin{column}{0.42\textwidth}
      \small The authority log already holds each operator's identity, takeover count, and takeover time.\\[0.25cm]
      {\normalsize Two minutes with a neighbor: what would ranking operators change?}
    \end{column}
  \end{columns}
  \mlsyscredit{Regulation (EU) 2024/1689 (AI Act), Annex III 4(b) and Art.\ 26(7), excerpts. Responsible Physical AI, \S17.7.3.}
  \note{[3 min] [CORE] NARRATE: the excerpt is real regulation text. Conformity assessment under the EU Machinery Regulation covers the machine's safety requirements for its declared use; it does not govern every later use of the machine's records. Scoring workers from them falls under the AI Act's employment category, with a duty to inform workers before use. ENGAGE: pairs, two minutes. Expected points: ranking turns safety records into performance scores; operators who expect to be ranked may take over less often, which removes protection and corrupts the evaluation record; it also needs its own legal classification. The chapter's line: a record declared for incident reconstruction may not be read to rank operators, and an access audit that finds such a read suspends that use and returns the premise to the adjudicator. A site that wants to evaluate workers runs a separate system with its own declaration and owner. WARN: the law named is European; other jurisdictions differ in detail, and the engineering premise does not. FLEX: [CORE].}
\end{frame}

\begin{frame}{Twenty machines create 17.3 person-hours of work a day}
  \begin{columns}[c]
    \begin{column}{0.56\textwidth}
      \centering\safeimg[width=\textwidth,height=4.4cm,keepaspectratio]{labor.pdf}\\[0.05cm]
      \raggedright\footnotesize Worker burden moves to rescue, watching, and sweeps: one assistant staffed 16 h, four 20 min sweeps (illustrative).
    \end{column}
    \begin{column}{0.42\textwidth}
      \centering\safeimg[width=\textwidth,height=3.0cm,keepaspectratio]{warehouse-drive-unit.jpg}\\[0.05cm]
      {\footnotesize A warehouse drive unit on display at an Amazon fulfilment center}
    \end{column}
  \end{columns}
  \mlsyscredit{Photo: Auledas, CC BY 4.0, Wikimedia Commons. Responsible Physical AI, Napkin Math 17.3 (inputs illustrative); hours derived, calcs.py.}
  \note{[2 min] [CORE] NARRATE: automating a task moves its labor. The assistant at 50 percent utilization spends 8 of 16 hours on requests and 8 watching and waiting; four sweeps of the restricted corridors add 1.3 h; total 17.3 person-hours a day, 2.2 full-time shifts, before maintenance, charging, or new demonstrations (calcs.py). The waiting hours are labor too: Bainbridge (1983) showed that people asked to watch and step in only for rare emergencies perform that job poorly. ENGAGE: the chapter reports, citing Burtch, Greenwood, and Ravindran (2025), warehouse robotics associated with fewer severe injuries and more non-severe ones, partly from a faster pace. What would the deployment need to log to test that here? (Expected: takeover, dispatch, and technician-hour records plus pace targets, and comparison sites no single deployment holds.) WARN: the photo is a different robot class; the labor pattern is the point. The labor block lists the work moved and the uses of worker data the deployment bars. FLEX: [CORE].}
\end{frame}

\begin{frame}{A decision with no owner refuses release}
  \begin{columns}[c]
    \begin{column}{0.6\textwidth}
      \scriptsize
      \begin{tabular}{@{}L{1.7cm}L{4.6cm}@{}}\toprule
        \textbf{Block} & \textbf{Holds, each with a named owner}\\\midrule
        Header & digests of manifest, authority, enforcement, fault records\\
        Decisions & type, owner, approval, evidence, expiry\\
        Human roles & time required and left, channels, verdict\\
        Permit & grantor, configuration, duties, disclosure log\\
        Exposure & groups, consent, per-group detection\\
        People & purpose, retention, readers, last audit\\
        Labor & hours moved, pace targets, barred uses\\\bottomrule
      \end{tabular}
    \end{column}
    \begin{column}{0.38\textwidth}
      \small The gate refuses on any of four findings:\\[0.1cm]
      \footnotesize\textcolor{figveto}{1.} a decision with no owner\\
      \textcolor{figveto}{2.} a cited role that is INFEASIBLE or UNKNOWN\\
      \textcolor{figveto}{3.} a relied-on group with no entry\\
      \textcolor{figveto}{4.} an expired or failed declaration\\[0.15cm]
      \centering\safeimg[width=0.8\textwidth,height=1.2cm,keepaspectratio]{thumb_record_chain.pdf}\\
      {\scriptsize\textcolor{midgray}{Record chain: after the release manifest}}
    \end{column}
  \end{columns}
  \mlsyscredit{Responsible Physical AI, \S17.9, Definition 17.3 and Table 2 (abridged). Record chain: Lecture 1; register: Lecture 18.}
  \note{[3 min] [CORE] NARRATE: the approval record is the signed record that names, for each decision behind the machine's admitted actions, the person who approved it and answers for it. It binds by digest the authority record, the enforcement record, the fault records, and the release manifest, then adds seven blocks, one per kind of decision in today's segments. The release gate reads it as a standing condition, so the four findings on the right refuse operation exactly as an unknown physical premise does, returning the premise to the adjudicator with its cause logged. A decision whose owner leaves the organization expires with the role. ENGAGE: how does it differ from the authority record? (Expected: the authority record shows who held each channel at each tick; the approval record shows who decided the machine could be in that situation at all.) LINK: entries that cannot close before release, a per-group sample still filling or a staffing plan waiting on a hire, leave with an operating restriction and enter the residual-claims register (Lecture 18) with an owner and a closure test. FLEX: [CORE].}
\end{frame}

% =============================================================================
% CLOSE
% =============================================================================
\begin{frame}{Which entries would have refused Cruise?}
  \mlsysactivity{Retrieval}{2 min}
  \begin{columns}[c]
    \begin{column}{0.34\textwidth}
      \centering\safeimg[width=\textwidth,height=2.6cm,keepaspectratio]{cruise-robotaxi.jpg}\\[0.05cm]
      {\footnotesize 6.1 m pullover; 15 days without a full account}
    \end{column}
    \begin{column}{0.64\textwidth}
      \only<1|handout:0>{\normalsize Close your notes. List the decisions behind the pullover and name the owner each needed.}%
      \only<2->{\footnotesize
      \begin{tabular}{@{}L{3.3cm}L{4.5cm}@{}}\toprule
        \textbf{Decision} & \textbf{What the gate checks}\\\midrule
        Post-collision response rule & a named owner, or refuse\\
        Operating domain & a per-group detection entry\\
        Commercial stage & grantor and date (CPUC, 10 Aug.\ 2023)\\
        Reporting duties & a disclosure-log owner and deadline\\\bottomrule
      \end{tabular}}
    \end{column}
  \end{columns}
  \mlsyscredit{Responsible Physical AI, \S17.1, \S17.4, and Checkpoint 17.3. Photo: Dietmar Rabich, CC BY-SA 4.0, Wikimedia Commons.}
  \note{[3 min] [CORE] Back to the hook. Step 1: retrieval, 60 seconds; Checkpoint 17.3 asks for the decisions behind the pullover and their owners. Step 2: one reasonable answer. The rule that a side impact triggers a pullover without checking beneath the vehicle needs a named owner; the public record does not show one, and a missing owner refuses. The operating domain needs a per-group detection entry. The commercial stage had a grantor and a date. The reporting duties existed in law; an owned disclosure log with deadlines is the part the regulators' penalties concerned. ENGAGE: write the first disclosure-log entry: who sends it, to whom, by when, with what evidence? (Expected: a named role; to the CPUC, the DMV, and NHTSA; within one day under NHTSA's order; the full video and the vehicle logs.) WARN: this is a teaching reconstruction from public sources, not a finding about any person. LINK: the hook asked who answered. The answer the approval record gives is the role named before the car moved. FLEX: [CORE].}
\end{frame}

\begin{frame}{Takeaways}
  \thumbline{thumb_decision_chain_step3.pdf}{Six decisions stand behind one admitted command; at 0.90 per link the chain is found only 53\% of the time, so each is recorded on purpose.}\\[0.12cm]
  \thumbline{thumb_queue_step2.pdf}{A takeover needing 4.5 s against a 0.85 s hazard adds blame and no protection; one assistant for 20 machines is busy half the time.}\\[0.12cm]
  \thumbline{thumb_ladder_step2.pdf}{Permission is granted in stages for a configuration, and withdrawn for unsafe motion or a withheld record.}\\[0.12cm]
  \thumbline{thumb_shortfall_step3.pdf}{Two frames late at 1.3 m/s is 43 mm, 0.42 m/s, and 32 J; resolving the gap takes about 1,059 people per group.}\\[0.12cm]
  \thumbline{thumb_retention.pdf}{Records about people declare purpose, retention, and readers; 20 machines make 17.3 person-hours of work a day.}
  \note{[1 min] [CORE] One line per segment. Ask which slide taught each line. Close on the fourth law: permission must be earned, and today it was earned by the decisions behind the evidence as well as the evidence itself. Energy delivered across the causal boundary cannot be recalled, so every owner is named before the first command is admitted.}
\end{frame}

\begin{frame}{Open questions}
  \begin{columns}[c]
    \begin{column}{0.62\textwidth}
      \small
      \begin{itemize}\setlength{\itemsep}{0.15cm}
        \item Who should grant the stages of a warehouse ladder when no public regulator does?
        \item Should a site choose its operating speed for its least-detected group, or keep that group out?
        \item When a machine's records could protect a worker in court and also score her, who decides which use wins?
      \end{itemize}
    \end{column}
    \begin{column}{0.36\textwidth}
      \centering\safeimg[width=\textwidth,height=3.4cm,keepaspectratio]{robot-welding-line.jpg}\\[0.05cm]
      {\footnotesize Robots spot-welding car bodies, BMW Leipzig, 2005}
    \end{column}
  \end{columns}
  \mlsyscredit{Photo: BMW Werk Leipzig, CC BY-SA 2.0 de, Wikimedia Commons.}
  \note{[1 min] [OPTIONAL] No single answer. The photo is the older answer to all three questions: keep people out of the cell. Physical AI brings the machine onto floors where people work and walk, so the decisions move from the fence to the approval record. The first question asks whether workers' representatives should hold a veto (\S17.4). The second sets speed against access (\S17.5 and \S17.6). The third is Nissenbaum's question of whether information flows as its context expects (\S17.7).}
\end{frame}

\begin{frame}{Read and explore}
  \small
  \begin{itemize}\setlength{\itemsep}{0.02cm}
    \item \textbf{Read:} Responsible Physical AI, \S17.1 to \S17.11; Checkpoints 17.1 to 17.3; Napkin Math 17.1 to 17.3 by hand.
    \item Thompson, ``The Problem of Many Hands,'' APSR 1980.
    \item Elish, ``Moral Crumple Zones,'' Engaging Science, Technology, and Society, 2019.
    \item Santoni de Sio and van den Hoven, ``Meaningful Human Control,'' Frontiers in Robotics and AI, 2018.
    \item Wilson, Hoffman, and Morgenstern, ``Predictive Inequity in Object Detection,'' 2019.
    \item Nissenbaum, \emph{Privacy in Context}, 2009.
  \end{itemize}
  \note{[0.5 min] [CORE] The chapter first, with its three checkpoints. Thompson names the problem of many hands; Elish shows case by case how blame settles on the person nearest the event; Santoni de Sio and van den Hoven state when human control is real; Wilson and colleagues measure per-group detection; Nissenbaum frames privacy as appropriate flow. FLEX: add Li et al. 2025 for a recent detection audit.}
\end{frame}

\begin{frame}[c]
  \begin{columns}[c]
    \begin{column}{0.36\textwidth}
      \centering\safeimg[width=\textwidth,height=2.4cm,keepaspectratio]{thumb_record_chain.pdf}
    \end{column}
    \begin{column}{0.60\textwidth}
      {\small\textcolor{lightgray}{NEXT LECTURE \textbar{} ENGINEERING PHYSICAL AI}}\\[0.5cm]
      {\LARGE\bfseries Why does a better model not, by itself, make a better machine?}
    \end{column}
  \end{columns}
  \note{[0.5 min] [CORE] Today's approval record leaves entries that cannot close before release: a per-group sample still filling, a staffing plan that depends on a hire, a crossing rule no sensor checks. Lecture 18 gives each open premise, physical or human, an owner, a restriction, and a closure test in the residual-claims register, and asks which one to retire next.}
\end{frame}

\end{document}
